\begin{definition}[Metric Space]\label{def:metric-space}
Let \(X\) be a nonempty set.  A function
\[
  d : X \times X \to \mathbb{R}_{\geq 0}
\]
is called a \emph{metric} on \(X\) if for all \(x,y,z \in X\),
\begin{enumerate}
  \item \(d(x,y) \geq 0\), and \(d(x,y) = 0\) if and only if \(x = y\)
        \hfill\textit{(positive definiteness)};
  \item \(d(x,y) = d(y,x)\)
        \hfill\textit{(symmetry)};
  \item \(d(x,y) \leq d(x,z) + d(z,y)\)
        \hfill\textit{(triangle inequality)}.
\end{enumerate}
A \emph{metric space} is a pair \((X,d)\), where \(X\) is a nonempty set and
\(d\) is a metric on \(X\).
\end{definition}

\begin{remark*}[Standard quantified statement]
\[
\begin{aligned}
\operatorname{Metric}(X,d)
\;\Longleftrightarrow\;
&X \neq \varnothing
\;\land\;
d:X\times X\to\mathbb{R}_{\geq 0}
\\
&{}\land
\forall x,y,z\in X{:}\;
d(x,y)\geq 0
\land
\bigl(d(x,y)=0 \Longleftrightarrow x=y\bigr)
\\
&{}\land d(x,y)=d(y,x)
\land d(x,y)\leq d(x,z)+d(z,y).
\end{aligned}
\]
\end{remark*}

\begin{remark*}[Interpretation]
A metric turns a set into a distance arena.  Positive definiteness says that
distance is never negative and that zero distance identifies exactly one
point.  Symmetry says that distance does not depend on direction.  The
triangle inequality says that moving through an intermediate point cannot
make the direct distance larger than the two-step route.
\end{remark*}

\begin{remark*}[Historical note]
This convention appears as Definition~2.2 in
Kainth's \emph{A Comprehensive Textbook on Metric Spaces}
\cite{KainthComprehensiveMetricSpaces2023}.
\end{remark*}

\DefinitionalRoot

\begin{definition}[Pseudo-Metric]\label{def:pseudo-metric}
Let \(X\) be a nonempty set.  A function
\[
  d : X \times X \to \mathbb{R}_{\geq 0}
\]
is called a \emph{pseudo-metric} on \(X\) if for all \(x,y,z \in X\),
\begin{enumerate}
  \item \(d(x,y) \geq 0\), and \(d(x,x) = 0\)
        \hfill\textit{(positive semi-definiteness)};
  \item \(d(x,y) = d(y,x)\)
        \hfill\textit{(symmetry)};
  \item \(d(x,y) \leq d(x,z) + d(z,y)\)
        \hfill\textit{(triangle inequality)}.
\end{enumerate}
\end{definition}

\begin{remark*}[Standard quantified statement]
\[
\begin{aligned}
\operatorname{PseudoMetric}(X,d)
\;\Longleftrightarrow\;
&X \neq \varnothing
\;\land\;
d:X\times X\to\mathbb{R}_{\geq 0}
\\
&{}\land
\forall x,y,z\in X{:}\;
d(x,y)\geq 0
\land d(x,x)=0
\\
&{}\land d(x,y)=d(y,x)
\land d(x,y)\leq d(x,z)+d(z,y).
\end{aligned}
\]
\end{remark*}

\begin{remark*}[Interpretation]
A pseudo-metric keeps the distance-like laws of nonnegativity, symmetry, and
the triangle inequality, but it weakens positive definiteness.  Distinct
points may have distance \(0\).  Thus a pseudo-metric measures separation only
up to possible zero-distance identifications.
\end{remark*}

\begin{remark*}[Historical note]
This convention appears as Definition~2.6 in
Kainth's \emph{A Comprehensive Textbook on Metric Spaces}
\cite{KainthComprehensiveMetricSpaces2023}.
\end{remark*}

\DefinitionalRoot

\begin{lemmabox}{Lemma}
\begin{lemma}[Metrics are Pseudo-Metrics]\label{lem:metrics-are-pseudo-metrics}
Every metric on a nonempty set \(X\) is a pseudo-metric on \(X\).
The converse is false: there exist pseudo-metrics that are not metrics.

\smallskip
\noindent
\hyperref[prf:metrics-are-pseudo-metrics]{\textit{Go to proof.}}
\end{lemma}
\end{lemmabox}

\begin{remark*}[Standard quantified statement]
\[
\begin{aligned}
&\forall X\,\forall d{:}\;
\operatorname{Metric}(X,d)
\Longrightarrow
\operatorname{PseudoMetric}(X,d),
\\
&\exists X\,\exists d{:}\;
\operatorname{PseudoMetric}(X,d)
\land
\neg\operatorname{Metric}(X,d).
\end{aligned}
\]
\end{remark*}

\begin{remark*}[Interpretation]
The implication is immediate because the metric axioms contain all
pseudo-metric axioms.  The only difference is the zero-distance law:
a metric forces \(d(x,y)=0\) exactly when \(x=y\), while a pseudo-metric only
requires \(d(x,x)=0\).
\end{remark*}

\begin{dependencies}
\begin{itemize}
  \item \hyperref[def:metric-space]{Definition: Metric Space}
  \item \hyperref[def:pseudo-metric]{Definition: Pseudo-Metric}
\end{itemize}
\end{dependencies}
